\documentclass[a4paper,12pt]{article}
\usepackage[T2A]{fontenc}
\usepackage[utf8]{inputenc}
\usepackage[russian]{babel}
\usepackage{amsmath,amssymb,mathtools}
\usepackage{geometry}
\geometry{left=22mm,right=22mm,top=23mm,bottom=25mm}
\usepackage{helvet}
\renewcommand{\familydefault}{\sfdefault}
\usepackage{microtype}
\usepackage{icomma}
\usepackage{booktabs,tabularx,array}
\usepackage{xcolor}
\usepackage{tikz}
\usepackage{hyperref}
\usepackage{fancyhdr}
\setlength{\headheight}{14pt}
\usepackage{enumitem}
\usepackage{parskip}
\definecolor{accent}{HTML}{005C64}
\definecolor{ink}{HTML}{263336}
\definecolor{error}{HTML}{B03939}
\definecolor{hint}{HTML}{235B9A}
\hypersetup{colorlinks=true,linkcolor=accent,urlcolor=accent}
\setlength{\parindent}{0pt}
\setlength{\parskip}{0.55em}
\setlist[itemize]{nosep,leftmargin=1.6em}
\pagestyle{fancy}
\fancyhf{}
\fancyhead[L]{\small\color{ink}Математика · Статистика}
\fancyhead[R]{\small\color{ink}04.11.2026}
\fancyfoot[C]{\footnotesize Лёвин Артём Александрович эксклюзивно для Николь Саркисьянц}
\fancyfoot[R]{\footnotesize\thepage}
\renewcommand{\headrulewidth}{0.3pt}
\renewcommand{\footrulewidth}{0pt}
\newcommand{\err}{\textcolor{error}{\textbf{Ошибка:}}}
\newcommand{\key}{\textcolor{hint}{\textbf{Ключевая идея:}}}
\newcommand{\res}{\textcolor{accent}{\textbf{Итог:}}}
\newcommand{\tasktitle}[1]{\phantomsection\label{task:#1}\section*{Задание №\,#1}}
\begin{document}
\thispagestyle{empty}
\begin{center}
\vspace*{54mm}
{\color{accent}\Large МАТЕМАТИКА}\par
\vspace{13mm}
{\Huge\bfseries Ревью домашней работы}\par
\vspace{12mm}
{\Large Описательная статистика}\par
\vspace{27mm}
{\large 4 ноября 2026 года}\par
\vspace{9mm}
{\normalsize Лёвин Артём Александрович}\par
\vfill
{\small Лёвин Артём Александрович эксклюзивно для Николь Саркисьянц}\par
\vspace{10mm}
\begin{tikzpicture}
 \draw[accent!65,line width=0.75pt] (-6.3,0) .. controls (-3.5,0.9) and (-2.7,-0.65) .. (0,0) .. controls (2.5,0.7) and (3.6,-0.8) .. (6.3,0);
\end{tikzpicture}
\end{center}
\newpage
\section*{Навигация по разборам}
В таблице представлены задания, в которых обнаружены ошибки. Номер задания ведёт к соответствующему разбору.

\vspace{6mm}
\footnotesize
\renewcommand{\arraystretch}{1.7}
\begin{tabularx}{\linewidth}{@{}c p{2.45cm} X p{2.75cm} p{3.05cm}@{}}
\toprule
№ & Тема & Причина ошибки & Ваш ответ & Правильный ответ\\
\midrule
\hyperref[task:5]{5} & Среднее арифметическое & Вместо пяти чисел взято шесть & $x=22$ & $x=10$\\
\hyperref[task:9]{9} & Дисперсия & Ошибка в квадрате отклонения; не учтены частоты & $D=7/6$; $\sigma\approx1{,}08$ & $D=17/9$; $\sigma\approx1{,}37$\\
\hyperref[task:10]{10} & Дисперсия и интервал & Неверное отклонение; не учтены частоты & $D\approx1{,}80$; $\sigma\approx1{,}34$ & $D=4{,}96$; $\sigma\approx2{,}23$\\
\bottomrule
\end{tabularx}
\normalsize
\vfill
\begin{center}\small\color{ink}В каждом разборе: источник ошибки, исправленный способ, проверка и одна задача на закрепление.\end{center}
\newpage
\thispagestyle{empty}
\vspace*{0.33\textheight}
\begin{center}
{\Huge\color{accent}\bfseries Разбор ошибок}
\end{center}
\newpage
\tasktitle{5}
\textbf{Условие.} Среднее арифметическое пяти чисел $5$, $11$, $x$, $15$, $19$ равно $12$. Найдите $x$.

\textbf{Ваше решение.} Записано:
\[
\frac{5+11+x+15+19}{\color{error}6}=12,
\qquad 5+11+x+15+19=72,
\qquad x=72-50=22.
\]
\textbf{Ваш ответ:} $x=22$.

\err\ В знаменателе записано $6$ вместо $5$.

\textbf{Где именно возникает ошибка.} Сумма чисел составлена верно: $5+11+x+15+19=50+x$. Первый неверный шаг --- деление суммы на $6$. В условии перечислено ровно пять чисел, включая неизвестное $x$.

\textbf{Почему это неверно.} Среднее арифметическое равно сумме \emph{всех} чисел, делённой на их \emph{количество}. Неизвестное число $x$ занимает одно место в наборе; наличие буквы не делает его дополнительным шестым числом. Из-за неверного знаменателя получилась сумма $72$ вместо $60$.

\key\ Перед составлением уравнения пересчитайте элементы набора и обозначьте их количество через $N$. Здесь $N=5$, поэтому
\[
\overline{x}=\frac{50+x}{5}.
\]

\textbf{Исправление.} Сохраняем верно найденную сумму $50+x$ и меняем только знаменатель:
\[
\frac{50+x}{5}=12.
\]

\textbf{Полное правильное решение.} Умножим обе части уравнения на $5$:
\[
50+x=12\cdot5=60.
\]
Вычтем $50$ из обеих частей:
\[
\boxed{x=10}.
\]

\textbf{Проверка.} Подставим $10$ в набор. Сумма пяти чисел равна
\[
5+11+10+15+19=60,
\qquad\frac{60}{5}=12.
\]
Получили среднее из условия.

\res\ $x=10$.

\subsection*{Задача на закрепление}
Среднее арифметическое пяти чисел $4$, $10$, $x$, $14$, $22$ равно $13$. Найдите $x$.
\newpage
\tasktitle{9}
\textbf{Условие.} В выборке число $1$ встречается два раза, число $3$ --- три раза, число $5$ --- один раз. По таблице частот вычислите среднее, дисперсию и стандартное отклонение.

\textbf{Ваше решение.} Верно составлена таблица частот $(2,3,1)$ и получено среднее $\overline{x}=16/6=8/3$. Затем найдены отклонения
\[
-1\frac23,\quad \frac13,\quad \frac73.
\]
Для первого отклонения записан квадрат $1\frac49$, а дисперсия указана как $D=\frac76$.

\textbf{Ваш ответ:} $D=\frac76$, $\sigma=\sqrt{\frac76}$.

\err\ Неверно возведено в квадрат первое отклонение; при вычислении дисперсии не учтено, сколько раз встречается каждое значение.

\textbf{Где именно возникает ошибка.} Последний бесспорно правильный шаг --- $1-\frac83=-\frac53=-1\frac23$. Первый неправильный шаг --- приписанный ему квадрат $1\frac49$. Действительно,
\[
\left(-1\frac23\right)^2
=\left(-\frac53\right)^2
=\frac{25}{9}=2\frac79.
\]
Кроме того, в записи дисперсии нужно учитывать частоты $2$, $3$ и $1$.

\textbf{Почему это неверно.} Чтобы возвести смешанное число в квадрат, сначала переводим его в неправильную дробь, затем возводим в квадрат числитель и знаменатель. При подсчёте дисперсии \emph{каждое наблюдение} вносит свой квадрат отклонения. Если значение повторяется три раза, его квадрат отклонения надо сложить трижды, то есть умножить на $3$.

\key\ Для таблицы значений $x_i$ и абсолютных частот $n_i$ действуют формулы
\[
N=\sum n_i,\qquad
\overline{x}=\frac{\sum n_i x_i}{N},\qquad
D=\frac{\sum n_i(x_i-\overline{x})^2}{N},
\qquad \sigma=\sqrt D.
\]

\textbf{Исправление.} Первое квадратное отклонение заменяем на $25/9$ и учитываем все частоты:
\[
D=\frac{2\cdot25/9+3\cdot1/9+1\cdot49/9}{6}.
\]

\textbf{Полное правильное решение.} Всего $N=2+3+1=6$ наблюдений. Среднее:
\[
\overline{x}=\frac{1\cdot2+3\cdot3+5\cdot1}{6}
=\frac{16}{6}=\frac83.
\]
\begin{center}
\renewcommand{\arraystretch}{1.3}
\begin{tabular}{@{}c c c c@{}}
\toprule
Значение $x_i$ & Частота $n_i$ & $x_i-\overline{x}$ & $(x_i-\overline{x})^2$\\
\midrule
$1$ & $2$ & $-5/3$ & $25/9$\\
$3$ & $3$ & $1/3$ & $1/9$\\
$5$ & $1$ & $7/3$ & $49/9$\\
\bottomrule
\end{tabular}
\end{center}
\[
D=\frac{2\cdot25+3\cdot1+1\cdot49}{9\cdot6}
=\frac{102}{54}=\boxed{\frac{17}{9}}.
\]
\[
\sigma=\sqrt{\frac{17}{9}}=\boxed{\frac{\sqrt{17}}3}\approx1{,}374.
\]

\textbf{Проверка.} Дисперсию можно пересчитать по формуле «среднее квадратов минус квадрат среднего»:
\[
\frac{2\cdot1^2+3\cdot3^2+1\cdot5^2}{6}
-\left(\frac83\right)^2
=9-\frac{64}{9}=\frac{17}{9}.
\]
Получилось то же значение.

\res\ $\overline{x}=\frac83$, $D=\frac{17}{9}$, $\sigma=\frac{\sqrt{17}}3$.

\subsection*{Задача на закрепление}
В выборке значение $1$ встречается два раза, значение $4$ --- три раза, значение $7$ --- один раз. Найдите среднее, дисперсию и стандартное отклонение.
\newpage
\tasktitle{10}
\textbf{Условие.} Значения $1$, $3$, $5$, $7$ встречаются с абсолютными частотами $3$, $2$, $3$, $2$. Найдите объём выборки, моду (или моды), медиану, среднее, размах, дисперсию, стандартное отклонение и интервал $[\overline{x}-\sigma;\,\overline{x}+\sigma]$.

\textbf{Ваше решение.} Найдены $N=10$, моды $1$ и $5$, медиана $4$, среднее $3{,}8$ и размах $6$. Далее при вычислении отклонений для значения $1$ записано $-2{,}38$; для дисперсии получено $1{,}79844$, стандартное отклонение округлено до $1{,}34$ и записан интервал $[2{,}46;\,5{,}14]$.

\textbf{Ваш ответ для ошибочной части:} $D=1{,}79844$, $\sigma\approx1{,}34$, интервал $[2{,}46;\,5{,}14]$.

\err\ В первом отклонении допущена арифметическая ошибка, а при подсчёте дисперсии квадраты отклонений сложены без умножения на частоты.

\textbf{Где именно возникает ошибка.} Последний правильный шаг --- $\overline{x}=3{,}8$. Первый неверный шаг --- разность $1-3{,}8=-2{,}38$. Правильно:
\[
1-3{,}8=\boxed{-2{,}8},\qquad (-2{,}8)^2=\boxed{7{,}84}.
\]
В тетради из неверной разности получено $5{,}6644$. Это значение тоже не соответствует правильному квадрату $7{,}84$.

\textbf{Почему это неверно.} При вычитании десятичных дробей нужно выровнять запятые: $1{,}0-3{,}8=-2{,}8$. Даже после исправления квадрата его недостаточно просто сложить с тремя остальными квадратами. Значение $1$ встречается трижды, $3$ --- дважды, $5$ --- трижды, $7$ --- дважды. Каждое повторение участвует в сумме дисперсии.

\key\ В таблице частот сначала найдите $N$ и $\overline{x}$. Затем для каждого значения вычислите отклонение, его квадрат и \emph{произведение квадрата на частоту}. Сумму последних произведений разделите на $N$.

\textbf{Исправление.} Сохраняем среднее $3{,}8$; пересчитываем первое отклонение и добавляем частоты в сумму:
\[
D=\frac{3\cdot7{,}84+2\cdot0{,}64+3\cdot1{,}44+2\cdot10{,}24}{10}.
\]

\textbf{Полное правильное решение.} Всего наблюдений
\[
N=3+2+3+2=10.
\]
Максимальная частота равна $3$, её имеют значения $1$ и $5$. Поэтому моды --- $1$ и $5$. Пятый элемент упорядоченной выборки равен $3$, шестой --- $5$, значит
\[
\operatorname{Me}=\frac{3+5}{2}=4.
\]
Среднее и размах:
\[
\overline{x}=\frac{1\cdot3+3\cdot2+5\cdot3+7\cdot2}{10}=\frac{38}{10}=3{,}8,
\qquad R=7-1=6.
\]
\begin{center}
\renewcommand{\arraystretch}{1.26}
\begin{tabular}{@{}c c c c c@{}}
\toprule
$x_i$ & $n_i$ & $x_i-3{,}8$ & $(x_i-3{,}8)^2$ & $n_i(x_i-3{,}8)^2$\\
\midrule
$1$ & $3$ & $-2{,}8$ & $7{,}84$ & $23{,}52$\\
$3$ & $2$ & $-0{,}8$ & $0{,}64$ & $1{,}28$\\
$5$ & $3$ & $1{,}2$ & $1{,}44$ & $4{,}32$\\
$7$ & $2$ & $3{,}2$ & $10{,}24$ & $20{,}48$\\
\bottomrule
\end{tabular}
\end{center}
\[
D=\frac{23{,}52+1{,}28+4{,}32+20{,}48}{10}
=\frac{49{,}60}{10}=\boxed{4{,}96}.
\]
\[
\sigma=\sqrt{4{,}96}=\frac{2\sqrt{31}}5\approx2{,}227.
\]
\[
[\overline{x}-\sigma;\,\overline{x}+\sigma]
\approx[3{,}8-2{,}227;\,3{,}8+2{,}227]
=\boxed{[1{,}57;\,6{,}03]}.
\]

\textbf{Проверка.} Среднее квадратов равно
\[
\frac{3\cdot1^2+2\cdot3^2+3\cdot5^2+2\cdot7^2}{10}
=\frac{194}{10}=19{,}4.
\]
Тогда $D=19{,}4-(3{,}8)^2=19{,}4-14{,}44=4{,}96$. Интервал симметричен относительно $3{,}8$.

\res\ $N=10$; моды $1$ и $5$; медиана $4$; среднее $3{,}8$; размах $6$; $D=4{,}96$; $\sigma\approx2{,}23$; интервал $[1{,}57;\,6{,}03]$.

\subsection*{Задача на закрепление}
Значения $0$, $2$, $4$, $6$ встречаются с частотами $2$, $3$, $2$, $1$ соответственно. Найдите объём выборки, моду, медиану, среднее, размах, дисперсию, стандартное отклонение и интервал $[\overline{x}-\sigma;\,\overline{x}+\sigma]$.
\end{document}
